% Part: first-order-logic
% Chapter: syntax-and-semantics
% Section: main-operator

\documentclass[../../../include/open-logic-section]{subfiles}

\begin{document}

\olfileid{fol}{syn}{mai}

\olsection{\printtoken{S}{main operator} لصيغة}

\begin{explain}
قد نحتاج في البحث إلى تعيين آخر مؤثر استُعمل في تأليف
!!a{formula}~$!A$؛ فهذا هو \emph{المؤثر الرئيس}
لـ~$!A$. وتقريبه إلى الذهن أنه المؤثر الذي يقع في ظاهر~$!A$
ويحيط بسائر تأليفها. فالمؤثر الرئيس لـ~$\lnot !A$ هو $\lnot$،
ولـ~$(!A \lor !B)$ هو $\lor$، وعلى ذلك يُقاس.
\end{explain}


\begin{defn}[!!^{main operator}]
\ollabel{def:main-op}
يُعرّف \emph{!!{main operator}} ل!!a{formula}~$!A$ كما يأتي:
\begin{enumerate}
\item \indcase*{!A}{!A}{ليس لـ~$\indfrm$ !!{main operator}.}

\tagitem{prvNot}{\indcase{!A}{\lnot !B}{!!{main operator} لـ~$\indfrm$
  هو~$\lnot$.}}{}

\tagitem{prvAnd}{\indcase{!A}{(!B \land !C)}{!!{main operator}
  لـ~$\indfrm$ هو~$\land$.}}{}

\tagitem{prvOr}{\indcase{!A}{(!B \lor !C)}{!!{main operator}
  لـ~$\indfrm$ هو~$\lor$.}}{}

\tagitem{prvIf}{\indcase{!A}{(!B \lif !C)}{!!{main operator}
  لـ~$\indfrm$ هو~$\lif$.}}{}

\tagitem{prvIff}{\indcase{!A}{(!B \liff !C)}{!!{main operator}
  لـ~$\indfrm$ هو~$\liff$.}}{}

\tagitem{prvAll}{\indcase{!A}{\lforall[x][!B]}{!!{main operator}
  لـ~$\indfrm$ هو~$\lforall$.}}{}

\tagitem{prvEx}{\indcase{!A}{\lexists[x][!B]}{!!{main operator}
  لـ~$\indfrm$ هو~$\lexists$.}}{}
\end{enumerate}
\end{defn}

والمراد في كل حالة \emph{الورود} المعيّن ل!!{main
operator}، لا مجرد نوع الرمز. فالصيغة $((!D \lif !E) \lif (!E \lif
!D))$ على صورة $(!B \lif !C)$؛ فيها $!B$ هي $(!D \lif !E)$ و$!C$
هي $(!E \lif !D)$. ومن ثم فالورود الثاني لـ~$\lif$ هو !!{main
operator}.

\begin{explain}
وبهذا عرّفنا تعريفًا \emph{عوديًا} دالةً تعيّن لكل !!{formula}
غير ذرية ورود !!{main operator} فيها. فإن التعريف الاستقرائي
لل!!{formula}s يقتضي أن تقع كل !!{formula}~$!A$ في إحدى حالات
\olref{def:main-op}؛ فثبت بذلك وجود !!a{main operator} لكل
!!{formula}~$!A$ غير ذرية. ثم إن كل !!{formula} لا تدخل إلا تحت
شرط واحد، وال!!{formula}s الأصغر التي تألّفت منها
$!A$ متعيّنة في كل حالة. فكان ورود !!{main operator}
لـ~$!A$ واحدًا لا غير؛ ولذلك صحّ أن ما عرّفناه دالة.
\end{explain}

ولل!!{formula}s أسماء يبيّنها
\olref{tab:main-op}، يختلف الاسم منها باختلاف الرمز الذي هو
!!{main operator} لها.\iftag{defNot,defOr,defAnd,defIf,defIff,defTrue,defFalse,defEx,defAll}
{ وأما المؤثرات المعرّفة، فلا تدخل في ال!!{formula}s دخولًا صوريًا؛
إنما هي اختصارات، فلا تكون على التحقيق مؤثرات رئيسة للصيغ.
ولهذا لا نفردها بحالات في البراهين التي تعمّ جميع
ال!!{formula}s.}

\begin{table}[!h]
\centering
\begin{tabular}{c | c | c}
!!^{main operator} & نوع ال!!{formula} & مثال\\
\hline
لا شيء & (!!{formula}) ذرية &
\iftag{prvFalse}{$\lfalse$،}{}
\iftag{prvTrue}{$\ltrue$،}{}
$\Atom{R}{t_1, \dots, t_n}$،
$\eq[{\OLId{latin-ordinary}{t}}_{\OLMathNumeral{1}}][{\OLId{latin-ordinary}{t}}_{\OLMathNumeral{2}}]$\\
$\lnot$ & نفي & $\lnot !A$ \\
$\land$ & عطف & $(!A \land !B)$ \\
$\lor$ & فصل & $(!A \lor !B)$ \\
$\lif$ & !!{conditional} & $(!A \lif !B)$ \\
$\liff$ & !!{biconditional} & $(!A \liff !B)$ \\
$\lforall[][]$ & (!!{formula}) كلية & $\lforall[{\OLId{latin-ordinary}{x}}][!A]$ \\
$\lexists[][]$ & (!!{formula}) وجودية & $\lexists[{\OLId{latin-ordinary}{x}}][!A]$
\end{tabular}
\caption{المؤثر الرئيس وأسماء ال!!{formula}s}
\ollabel{tab:main-op}
\end{table}

\end{document}
